\documentclass[aspectratio=169]{beamer}

% The 'handout' option collapses all overlays into single slides
%\documentclass[aspectratio=169, handout]{beamer}

\usepackage{amsmath, amssymb, amsthm, mathtools}
\usepackage{algorithm, algorithmic}
\usepackage{tikz}
\usetikzlibrary{arrows.meta,positioning,fit,calc,shapes}

\newtheorem{proposition}[theorem]{Proposition}

% --- Standard Beamer Theme ---
\usetheme{Madrid}
\usecolortheme{default}

\include{macros} % Assuming you have this file


\title[LinUCB]{Linear Bandits}
%\subtitle{Part 1: The Assumption Ladder \& Bellman Completeness}
\author[Sham \& Kiant\'{e}]{Kiant\'{e} Brantley \& Sham Kakade}
\date{}


\begin{document}

\begin{frame}
\titlepage  
\end{frame}

\begin{frame}[noframenumbering]
  \frametitle{Outline}
  \tableofcontents
\end{frame}


\section{Recap}

\begin{frame}\frametitle{Generalization in RL}
  
\begin{itemize}
\item (distribution free) Agnostic learning is not possible in RL:\\
we showed that to get $O(\log |\Pi|)$ sample complexity we need either:
\begin{itemize}
\item $\textrm{poly}(|\Scal|)$ samples OR
\item $\textrm{exp}(H)$ samples.
\end{itemize}
in order to learn the best policy in some policy class.
\item  upshot: we need stronger  assumptions for RL analysis.
\end{itemize}

\end{frame}


\section{Linear Bandits}

\begin{frame}[noframenumbering]
  \frametitle{Outline}
  \tableofcontents[currentsection]
\end{frame}

\subsection{Setting}

\begin{frame}\frametitle{Handling Large Actions Spaces}

\begin{itemize}
\item On each round, we must choose a decision $x_t \in D\subset R^d$. 
\item \pause Obtain a reward $r_t \in [-1, 1]$, where
%for all $x \in D$,
\alert{\[
\E[r_t|x_t=x] = \mu^\star \cdot x \in [-1, 1],
\]}
 \pause 
\begin{itemize}
\item so the the conditional expectation of $r_t$ is linear)
\item Also, we have the \emph{noise sequence},
\[
\eta_t = r_t - \mu^\star \cdot x_t  
\]
is i.i.d noise.
\end{itemize}
\end{itemize}
{\small model due to Abe \& Long '99}
\end{frame}

\begin{frame}\frametitle{Our Objective}

If $x_0, \ldots x_{T-1}$ are our decisions, then 
our \alert{cumulative regret} is
\[
R_T = T \mu^\star\cdot x^\star - \sum_{t=0}^{T-1} \mu^\star\cdot x_t
%R_T = \mu^\star\cdot x^\star - \sum_{t=0}^{T-1} r_t
\]
where $x^\star \in D$ is an optimal decision for $\mu^\star$, i.e.
\[
x^\star \in \argmax_{x\in D} \mu^\star \cdot x
\]
\end{frame}


\subsection{LinUCB}

\begin{frame}[noframenumbering]
  \frametitle{Outline}
  \tableofcontents[currentsubsection]
\end{frame}


\begin{frame}
\frametitle{The ``Confidence Ball''}

After $t$ rounds, define our uncertainty region 
$\textsc{Ball}_t$: with center, $\widehat{\mu}_t$, and shape,
$\Sigma_t$, using the $\lambda$-regularized least squares solution:
\begin{align*}%\label{eq:regression}
\widehat{\mu}_t &= \arg\min_{\mu} \sum_{\tau = 0}^{t-1}
\| \mu \cdot x_\tau - r_\tau \|_2^2 + \lambda\|\mu\|_2^2\\
&= \Sigma_t^{-1}\sum_{\tau = 0}^{t-1}r_\tau x_\tau,\\
\Sigma_t &= \lambda I + \sum_{\tau = 0}^{t-1}x_\tau x_\tau^\top,
\, \mathrm{with } \, \, \, \Sigma_0 = \lambda I.
\end{align*}

\pause
Define the uncertainty region:
\begin{align*}
\textsc{Ball}_{t} &= \left\{ \mu \mid
(\widehat \mu_t - \mu)^\top\Sigma_t(\widehat \mu_t - \mu)
\leq \beta_t \right\},  
%\textsc{Ball}_{t} = \left\{ 
%(\widehat \mu_t - \mu^\star)^\top\Sigma_t^{-1}(\widehat \mu_t - \mu^\star)
%  \leq \beta_t \right\}, 
\end{align*}
where $\beta_t$ is a parameter of the algorithm.

\end{frame}

\begin{frame}
\frametitle{LinUCB (the algo)}

\begin{enumerate}
\item Input: $\lambda$, $\beta_t$
  \item For $t=0,1, \ldots$
    \begin{enumerate} 
  \item Execute 
    \[
    x_t = \argmax_{x\in D}\max_{\mu\in \textsc{Ball}_t} \mu\cdot x 
    \]
    and observe the reward $r_t$.
  \item Update $\textsc{Ball}_{t+1}$.
  \end{enumerate}
\end{enumerate}

\end{frame}

\subsection{An Optimal Regret Bound}

\begin{frame}\frametitle{LinUCB Regret Bound}


Sublinear regret: $R_T \leq O^\star(d \sqrt{T})$ \\
\alert{poly dependence on $d$} , no dependence on the cardinality $|D|$.

\pause
\begin{theorem}~\label{theorem:sqrtregret}
Suppose: $|\mu^\star\cdot x| \leq 1$ and $\|x\|\leq B$ for all $x\in
D$; that the noise is $\sigma^2$ sub-Gaussian; and that
$\|\mu^\star\|\leq W$. Set
$\lambda =\sigma^2/W^2$ and
\[
\beta_t := \sigma^2\Big(2
+ 4 d\log\left( 1 + \frac{tB^2W^2}{ d} \right)
+8   \log(4/ \delta)\Big).
\]
With probability greater than $1-\delta$, that for all $T\geq 0$, 
\[
R_T \leq c \sigma \sqrt{T} \left( d\log\left( 1 + \frac{TB^2W^2}{ d\sigma^2} \right)
+   \log(4/ \delta) \right)
\]
where $c$ is an absolute constant. 
\end{theorem}
{\small due to Dani, Hayes, K. '09 }

\end{frame}



\section{Analysis}

 \begin{frame}[noframenumbering]
    \frametitle{Outline}
    \tableofcontents[currentsection]
  \end{frame}


\begin{frame}\frametitle{Confidence}

In establishing the upper bounds there are two main propositions from
which the upper bounds follow. The first is in showing that the
confidence region is valid. 

\begin{proposition} \label{proposition:confidence} 
(Confidence) Let $\delta>0$.  We have that
%\item For ConfidenceBall$_2$,
\[
\Pr(\forall t, \, \mu^\star \in \textsc{Ball}_t ) \ge 1 - \delta.
\]
\end{proposition}

\end{frame}

\begin{frame}\frametitle{Sum of Squares Regret Bound}


Assuming the confidence event holds,
the following controls on the growth of the regret.  



\begin{proposition}\label{proposition:l2regretbd}
(Sum of Squares Regret Bound) Define:
\[
\mathrm{regret}_t = \mu^\star \cdot x^{*} - \mu^\star \cdot x_t 
\]
Suppose $\|x\|\leq B$ for $x\in
D$. Suppose $\beta_t$ is increasing and larger than $1$.
Suppose $\mu^\star \in \textsc{Ball}_t$ for all $t$, then 
\[
\sum_{t=0}^{T-1} \mathrm{regret}_t^2  \le %4n \beta_T \ln T  
4\beta_T d\log\left( 1 + \frac{TB^2}{ d \lambda } \right)
\]
\end{proposition}

\end{frame}

\begin{frame}\frametitle{Completing the Proof}
 
\begin{proof}[Proof of Theorem~\ref{theorem:sqrtregret}]
With the two previous Propositions,  along with
the Cauchy-Schwarz 
inequality, we have, with probability at least $1-\delta$,
\begin{align*}
&R_T = \sum_{t=0}^{T-1} \mathrm{regret}_t
\le \sqrt{ T \sum_{t=0}^{T-1} \mathrm{regret}_t^2 }
\le \sqrt{4T\beta_T d\log\left( 1 + \frac{TB^2}{ d \lambda } \right)}.
\end{align*}
The remainder of the proof follows from using our chosen value of
$\beta_T$ and algebraic manipulations.
\end{proof}

\end{frame}

\subsection{Regret Analysis}
 \begin{frame}[noframenumbering]
    \frametitle{Outline}
    \tableofcontents[currentsubsection]
  \end{frame}

\begin{frame}\frametitle{``Width'' of Confidence Ball}

\begin{lemma} \label{lemma:width}
Let $x \in D$. If $\mu \in \textsc{Ball}_t$ and $x \in D$. Then 
\[
|(\mu - \muhat_t)^\top x| \le \sqrt{\beta_t x^\top \Sigma_t^{-1}x}
\]
\end{lemma}

\pause
\begin{proof}  
By
Cauchy-Schwarz, we have:
\begin{align*}
&|(\mu - \muhat_t)^\top x| = 
|(\mu - \muhat_t)^\top  \Sigma_t^{1/2}\Sigma_t^{-1/2}x| 
= |(\Sigma_t^{1/2}(\mu - \muhat_t))^\top \Sigma_t^{-1/2}x| \\
& \le \| \Sigma_t^{1/2}(\mu - \muhat_t) \| \|\Sigma_t^{-1/2}x\| 
= \|\Sigma_t^{1/2}(\mu - \muhat_t)\| 
\sqrt{x^\top  \Sigma_t^{-1} x} 
\le \sqrt{\beta_t x^\top \Sigma_t^{-1}x} 
\end{align*}
where the last inequality holds since $\mu \in \textsc{Ball}_t$. 
\end{proof}


\end{frame}


\begin{frame}\frametitle{Instantaneous Regret Lemma}

Define
\[
w_t := \sqrt{x_t^\top  \Sigma_t^{-1} x_t}
\]
which is the ``normalized width'' at time $t$ in the direction of our decision.
%We now see that the width, $2 \sqrt{\beta_t} w_t$, 
%is an upper bound for the instantaneous regret.

\pause
\begin{lemma} \label{lem:regretvswidth}
Fix $t\le T$. If $\mu^\star  \in \textsc{Ball}_t$, then
\[
\mathrm{regret}_t \le 2\min{(\sqrt{\beta_t} w_t , 1)}
\le 2 \sqrt{\beta_T}\min{(w_t , 1)}
\]
\end{lemma}

\pause
\begin{proof}
Let $\mutilde\in \textsc{Ball}_t$ denote the vector which  maximizes
the dot product $\mutilde^\top  x_t$.  By the choice of
$x_t$, we have
\[
\mutilde^\top x_t 
=\max_{\mu \in \textsc{Ball}_t} \mu^\top  x_t
= \max_{x\in D}\max_{\mu \in \textsc{Ball}_t} \mu^\top  x
\ge (\mu^\star)^\top  x^{*},
%\mutilde^\top x_t = \max_{\mu \in \textsc{Ball}_t} \max_{x\in D} \mu^\top  x
%\ge (\mu^\star)^\top  x^{*},
\]
where the inequality used the hypothesis $\mu^\star  \in \textsc{Ball}_t$.
Hence, 
\begin{align*}
&\mathrm{regret}_t =  (\mu^\star)^\top  x^{*} -(\mu^\star)^\top  x_t
\le ( \mutilde - \mu^\star   )^\top  x_t \\
&\quad \quad =  (\mutilde-\muhat_t )^\top x_t + (\muhat_t-\mu^\star)^\top x_t 
\le 2\sqrt{\beta_t} w_t
\end{align*}
where the last step follows from Lemma~\ref{lemma:width}
since $\mutilde$ and $\mu^\star $ are in $\textsc{Ball}_t$. Since $r_t \in [-1, 1]$,
$\mathrm{regret}_t$ is always at most $2$ and the first inequality follows. 
The final inequality is due to that $\beta_t$ is increasing and larger than $1$.
\end{proof}

\end{frame}


\begin{frame}\frametitle{Geometric Argument: Part 1}

The next two lemmas give us 'geometric' potential function argument, where can bound the sum of widths independently of
the choices made by the algorithm. 

\pause
\begin{lemma} \label{lem:volumechange}
We have:
\[
\det \Sigma_{T} = \det \Sigma_{0}\prod_{t = 0}^{T-1} (1 + w_t^2).
\]
\end{lemma}

\pause
\begin{proof}
By the definition of $\Sigma_{t+1}$, we have
\begin{align*}
&\det \Sigma_{t+1}  = \det (\Sigma_t + x_t x_t^\top )
 = \det ( \Sigma_t^{1/2} (I + \Sigma_t^{-1/2}x_t x_t^\top \Sigma_t^{-1/2}) \Sigma_t^{1/2}) \\
& \quad = \det(\Sigma_t) \det (I + \Sigma_t^{-1/2}x_t (\Sigma_t^{-1/2}x_t)^\top ) 
 = \det(\Sigma_t) \det(I + v_t v_t^\top ),
\end{align*}
where $v_t := \Sigma_t^{-1/2}x_t$. 
Now observe that $v_t^\top v_t = w_t^2$ and $\ldots$
\iffalse
\[
(I  + v_t v_t^\top  ) v_t = v_t + v_t(v_t^\top v_t) = (1+ w_t^2) v_t
\]
Hence $(1+ w_t^2)$ is an eigenvalue of $I  + v_t v_t^\top $. Since
$v_tv_t^\top $ is a rank one matrix, all other eigenvalues 
of $I + v_t v_t^\top $ equal 1.
Hence, $\det (I  + v_t v_t^\top  )$ is $(1+ w_t^2)$, 
implying
$\det \Sigma_{t+1} = (1 + w_t^2) \det \Sigma_t$.
%Recalling that $\Sigma_{1}$ is the identity matrix, 
The result follows by induction.
\fi
\end{proof}

\end{frame}


\begin{frame}\frametitle{Geometric Argument: Part 2}
\vspace*{-0.1cm}
\begin{lemma}\label{lem:detA}
For any sequence $x_0, \ldots x_{T-1}$ such that, for $t<T$,
$\|x_t\|_2 \leq B$, we have: 
\begin{align*}
\log\Big(\det \Sigma_{T-1}/\det \Sigma_{0}\Big) = \log\det\left( I + \frac{1}{\lambda} \sum_{t=0}^{T-1}x_t x_t^{\top}\right)
\leq  d\log\left( 1 + \frac{TB^2}{ d \lambda } \right).
\end{align*}
\end{lemma}
\vspace*{-0.1cm}
\pause
{\small 
\begin{proof}
Denote the eigenvalues of $ \sum_{t=0}^{T-1}x_t x_t^{\top}$ as
$\sigma_1, \dots \sigma_d$, and note:
\begin{align*}
\sum_{i=1}^d \sigma_i = \mathrm{Trace}\Big( \sum_{t=0}^{T-1}x_t x_t^{\top}  \Big) 
= \sum_{t=0}^{T-1}\|x_t\|^2 \leq T B^2.
\end{align*} 
Using the AM-GM inequality,
\begin{align*}
&\log\det\Big( I + \frac{1}{\lambda} \sum_{t=0}^{T-1}x_t x_t^{\top}\Big) = \log\Big( \prod_{i=1}^d \left(1 + \sigma_i / \lambda \right) \Big) \\
&= d\log\Big( \prod_{i=1}^d \left(1 + \sigma_i / \lambda \Big)
  \right)^{1/d} 
\leq d\log\Big( \frac{1}{d}\sum_{i=1}^d \left(1 + \sigma_i / \lambda \right) \Big) 
\leq d \log\Big( 1 + \frac{ TB^2  }{d\lambda}  \Big)
\end{align*} 
which concludes the proof. 
\end{proof}
}
\end{frame}

\begin{frame}\frametitle{Proving ``sum of squares
    regret'' Proposition}

\begin{proof}[Proof of Proposition~\ref{proposition:l2regretbd}]
Assume $\mu^\star  \in \textsc{Ball}_t$ for all $t$. We have:
\begin{align*}
&\sum_{t=0}^{T-1} \mathrm{regret}_t^2 
\le \sum_{t=0}^{T-1} 4\beta_t \min (w_t^2, 1)
\le 4\beta_T \sum_{t=0}^{T-1}  \min (w_t^2, 1)\\
&\quad \le 8\beta_T \sum_{t=0}^{T-1}  \ln(1+w_t^2) 
\le 8\beta_T \log\Big(\det \Sigma_{T-1}/\det \Sigma_{0}\Big)\\
&\quad=8\beta_T d\log\left( 1 + \frac{TB^2}{ d \lambda } \right)
\end{align*}
where the first inequality follow from by
Lemma~\ref{lem:regretvswidth}; the second from that $\beta_t$ is an
increasing function of $t$; the third uses that for $0 \le y \le 1$,
$\ln(1 + y) \ge y/2$; the final two inequalities follow by
Lemmas~\ref{lem:volumechange} and~\ref{lem:detA}. 
\end{proof}
\end{frame}


\subsection{Confidence Analysis}
 \begin{frame}[noframenumbering]
    \frametitle{Outline}
    \tableofcontents[currentsubsection]
  \end{frame}



\begin{frame}\frametitle{Confidence [Proof of Proposition~\ref{proposition:confidence}] }

\begin{proof}
Since $r_\tau=x_\tau \cdot\mu^\star +\eta_\tau $, we have:
\begin{align*}
&\widehat \mu_t - \mu^\star 
= \Sigma_t^{-1}\sum_{\tau = 0}^{t-1}r_\tau x_\tau- \mu^\star 
= \Sigma_t^{-1}\sum_{\tau = 0}^{t-1} x_\tau (x_\tau \cdot\mu^\star +\eta_\tau)- \mu^\star \\
&= \Sigma_t^{-1}\left( \sum_{\tau = 0}^{t-1} x_\tau (x_\tau)^\top  \right) \mu^\star 
- \mu^\star 
+\Sigma_t^{-1}\sum_{\tau = 0}^{t-1} \eta_\tau x_\tau\\
&= \lambda \Sigma_t^{-1}\mu^\star 
+\Sigma_t^{-1}\sum_{\tau = 0}^{t-1} \eta_\tau x_\tau
\end{align*}

\pause
By the triangle inequality, 
\begin{align*}
% \left\| \widehat \mu_t - \mu^\star    \right\|_{\Sigma_t^{-1}}
\sqrt{(\widehat \mu_t - \mu^\star)^\top\Sigma_t(\widehat \mu_t - \mu^\star)}
&\leq \left\| \lambda \Sigma_t^{-1/2}\mu^\star  \right\|
+ \left\| \Sigma_t^{-1/2}\sum_{\tau = 0}^{t-1} \eta_\tau x_\tau \right\| \\
& \leq \sqrt{\lambda}\|\mu^\star\| 
\quad\quad \quad + \quad\quad \quad??.
%+ \sqrt{2 \sigma^2\log\left(\det(\Sigma_t)\det(\Sigma^0)^{-1} / \delta_t \right) }.
\end{align*} 
How can we bound ``??''
To be continued...
\end{proof}
\end{frame}

\begin{frame}\frametitle{Self-Normalizing Sum}

\begin{lemma}[Self-Normalized Bound for Vector-Valued Martingales]
(Abassi et. al '11)
Suppose $\{\varepsilon_i\}_{i=1}^{\infty}$ are mean zero random
variables (can be generalized to martingales), and $\varepsilon_i$ is
bounded by $\sigma$. Let
$\{X_i\}_{i=1}^{\infty}$ be a stochastic process. 
Define $\Sigma_t = \Sigma_0 + \sum_{i=1}^{t} X_iX_i^{\top}$. With probability
at least $1-\delta$, we have for all $t\geq 1$: 
\begin{align*}
\left\|  \sum_{i=1}^{t} X_i \varepsilon_i \right\|^2_{\Sigma_t^{-1}} \leq \sigma^2 \log\left(\frac{ \det( \Sigma_t) \det(\Sigma_0)^{-1} }{ \delta^2  }\right).
\end{align*}
\end{lemma}

\end{frame}


\begin{frame}\frametitle{Continued... [Proof of Proposition~\ref{proposition:confidence}] }

\begin{proof}
\begin{align*}
(\widehat \mu_t - \mu^\star)^\top \Sigma_t (\widehat \mu_t - \mu^\star)
&\leq \left\| \lambda \Sigma_t^{-1/2}\mu^\star  \right\|
+ \left\| \Sigma_t^{-1/2}\sum_{\tau = 0}^{t-1} \eta_\tau x_\tau \right\| \\
& \leq \sqrt{\lambda}\|\mu^\star\| 
+ \sqrt{2 \sigma^2\log\left(\det(\Sigma_t)\det(\Sigma^0)^{-1} / \delta_t \right) }.
\end{align*} 

We seek to lower bound $\Pr(\forall t, \, \mu^\star \in
\textsc{Ball}_t )$. Assign failure probability $\delta_t=(3/\pi^2)/t^2$ for the
$t$-th event, which gives us:
\begin{align*}
&1-\Pr(\forall t, \, \mu^\star \in \textsc{Ball}_t ) 
= \Pr(\exists t, \, \mu^\star \notin \textsc{Ball}_t )
\leq \sum_{t=1}^{\infty} \mathrm{Pr}( \mu^\star \notin \textsc{Ball}_t) \\
&\quad <\sum_{t=1}^{\infty} (1/t^2)
(3/\pi^2) = 1/2.
\end{align*}
This along with Lemma~\ref{lem:detA}  completes the proof.
\end{proof}
\end{frame}


\end{document}

